%% chapters/02-writing-content.tex
\chapter{Writing Content}\label{chap:writing-content}

This chapter shows the building blocks you will use most often: headings,
text, lists, callout boxes and theorem-like environments.

\section{Headings}\label{sec:headings}

The document is organised with the following commands, from largest to
smallest: \cmd{chapter}, \cmd{section}, \cmd{subsection} and
\cmd{subsubsection}. Chapters, sections and subsections are numbered and
listed in the table of contents; deeper levels are not, unless you change the
depth (see \cref{app:customization}).

\subsection{A Subsection}

Subsections group related paragraphs inside a section.

\subsubsection{A Subsubsection}

This level is not numbered by default. Use \cmd{paragraph} for a run-in
heading at an even lower level.

\section{Text and Emphasis}\label{sec:text}

Use \cmd{emph} for emphasis and \cmd{textbf} for strong emphasis. Quotation
marks should be typed with \cmd{enquote}, which follows the conventions of the
document language: \enquote{a quotation with \enquote{a nested one}}. Insert a
non-breaking space with a tilde, as in Figure~1, so that a label and its
number are never split across lines. Footnotes appear at the bottom of the
column.\footnote{Like this one.}

\section{Lists}\label{sec:lists}

Bulleted, numbered and description lists share the same compact style:

\begin{itemize}
  \item A bulleted item.
  \item Another bulleted item.
\end{itemize}

\begin{enumerate}
  \item A numbered item.
  \item Another numbered item.
\end{enumerate}

\begin{description}
  \item[Term] A short explanation of the term.
  \item[Another term] A second explanation.
\end{description}

\section{Callout Boxes}\label{sec:callouts}

Callout boxes draw attention to information that deserves to stand out. Three
kinds are available.

\begin{note}
Use a \emph{note} for supplementary information the reader may find useful.
\end{note}

\begin{tip}
Use a \emph{tip} for a recommendation or a shortcut.
\end{tip}

\begin{warning}
Use a \emph{warning} for a pitfall that could cause errors or lost work.
\end{warning}

\noindent Each box is an environment named after its kind:

\begin{lstlisting}[language={[LaTeX]TeX}]
\begin{note}
Your text here.
\end{note}
\end{lstlisting}

\section{Definitions, Examples and Exercises}\label{sec:theorems}

Four theorem-like environments are provided. They share one counter that
restarts in every chapter, so \cref{def:template} is followed by
\cref{ex:template} and \cref{thm:square}.

\begin{definition}[Template]\label{def:template}
A \emph{template} is a reusable document skeleton that fixes the layout and
leaves the content to the author.
\end{definition}

\begin{example}\label{ex:template}
This book is an example of a template: the design is fixed, and only the
text, figures and references change from one book to the next.
\end{example}

\begin{theorem}\label{thm:square}
For all real numbers $a$ and $b$, $(a+b)^2 = a^2 + 2ab + b^2$.
\end{theorem}

\begin{proof}
Expand the product $(a+b)(a+b)$ and collect the terms.
\end{proof}

\begin{exercise}\label{ex:practice}
Add a new chapter to this book (\cref{sec:adding-chapter}) and reference
\cref{thm:square} from it.
\end{exercise}

\section{Columns and Margin Notes}\label{sec:layout}

The book is typeset in two columns by default, which keeps lines short and
readable on a landscape page. To switch to a single column, delete the
\texttt{twocolumn} option in \file{main.tex}. The page then has one text block
and a wide outer margin, where you can place short notes with
\cmd{marginline}. In two-column mode the margins are narrow, so margin notes
are not recommended.

\begin{note}
Wide material such as a large table or figure should use the starred
environments, \texttt{table*} and \texttt{figure*}, which span both columns
(see \cref{chap:figures}).
\end{note}
